\documentclass[11pt]{article}
\usepackage[a4paper,margin=20mm]{geometry}
\usepackage{fontspec}
\setmainfont{Noto Serif CJK SC}
\XeTeXlinebreaklocale "zh"
\XeTeXlinebreakskip=0pt plus 1pt
\usepackage{amsmath,amssymb,amsthm,mathtools,tikz-cd,enumerate,paralist,xurl,hyperref}
\setlength{\emergencystretch}{4em}
\newcommand{\identity}{\mathrm{id}}
\newcommand{\Obj}{\operatorname{Ob}}
\newcommand{\Mor}{\operatorname{Mor}}
\newcommand{\Hom}{\operatorname{Hom}}
\newcommand{\End}{\operatorname{End}}
\newcommand{\Ker}{\operatorname{ker}}
\newcommand{\Coker}{\operatorname{coker}}
\newcommand{\opp}{\mathrm{op}}
\newcommand{\munit}{\mathbf{1}}
\newcommand{\Z}{\mathbb{Z}}
\newcommand{\cate}[1]{\mathbf{#1}}
\newcommand{\cated}[1]{\mathbf{#1}\text{-}}
\newcommand{\dcate}[1]{\text{-}\mathbf{#1}}
\newcommand{\dotimes}[1]{\mathbin{\mathop{\otimes}\limits_{#1}}}
\newcommand{\rightiso}{\xrightarrow{\sim}}
\renewcommand{\index}[2][]{\relax}
\newcommand{\CorpusNumber}{0}
\newtheorem{theorem}{定理}
\newtheorem{lemma}[theorem]{引理}
\newtheorem{proposition}[theorem]{命题}
\newtheorem{corollary}[theorem]{推论}
\theoremstyle{definition}
\newtheorem{definition}[theorem]{定义}
\newtheorem{definition-proposition}[theorem]{定义--命题}
\newtheorem{example}[theorem]{例}
\newtheorem{remark}[theorem]{注记}
\renewcommand{\thetheorem}{\CorpusNumber}
\renewcommand{\proofname}{证明}
\begin{document}
\section*{293: Beck monadicity}
\noindent\textbf{作者：} 李文威 (Wen-Wei Li). \textbf{作品：}《代数学方法》第二卷，第七章。
\par\noindent\textbf{来源与版本：} \url{https://github.com/wenweili/AlJabr-2/blob/e03224e2e867a49dc60c7975d890079068a30466/chapter7.tex#L1574}, commit \texttt{e03224e2e867a49dc60c7975d890079068a30466}.
\par\noindent\textbf{许可：} Creative Commons Attribution 4.0 International; copyright 2024 李文威. \url{https://creativecommons.org/licenses/by/4.0/}. 作者不保证无误；见附带 LICENSE。
\par\noindent\textbf{整理说明（非作者正文）：} 下列题面、证明及局部支撑正文来自作者原生 TeX；保留原语言、顺序及数学内容。仅改独立排版、原编号引用与索引命令；未生成或补写证明。所有文字与数学图表均可编辑。标准先修仅引用，不递归重证。
\par\noindent\textbf{难度说明（整理者）：} H2: The proof constructs a quasi-inverse to the comparison functor using split coequalizers and proves the required natural isomorphisms. This reconstructs an entire category of algebras; it is not a routine verification of a given adjunction.
\subsection*{原命题与完整人类证明}
\def\CorpusNumber{7.6.14}
\begin{theorem}[J.\ M.\ Beck]\label{prop:Beck}
	\index{Beck dingli@Beck 定理}
	设函子 $G: \mathcal{D} \to \mathcal{C}$ 有左伴随 $F$. 以下陈述等价:
	\begin{enumerate}[(i)]
		\item 伴随对 $(F, G)$ 是单子的 (定义 7.6.8);
		\item $G$ 对所有 $G$-分裂对 $f, g: X \rightrightarrows Y$ 都生相应的余等化子;
		\item $G$ 是保守的 (定义 1.5.4), 所有 $G$-分裂对 $f, g: X \rightrightarrows Y$ 在 $\mathcal{D}$ 中都有余等化子 $X \rightrightarrows Y \to Z$, 而且后者在 $G$ 之下的像给出 $(Gf, Gg)$ 的余等化子.
	\end{enumerate}

	当以上任一条件成立时, $\mathbb{K}: \mathcal{D} \to \mathcal{C}^T$ 的一个拟逆函子 $\mathbb{L}$ 可以按以下方式描述: 设 $(M, a) \in \Obj(\mathcal{C}^T)$, 则有余等化子图表
\renewcommand{\theequation}{7.6.1}
	\begin{equation}\label{eqn:Beck-aux1}\begin{tikzcd}
		FGF(M) \arrow[shift left, r, "Fa"] \arrow[shift right, r, "\varepsilon_{FM}"'] & F(M) \arrow[r] & \mathbb{L}(M, a).
	\end{tikzcd}\end{equation}

	对于由 $(F, G)$ 确定的余单子, 判准完全是对偶的.
\end{theorem}
\begin{proof}
	我们先证明 (i) $\implies$ (ii). 设 $\mathbb{K}: \mathcal{D} \to \mathcal{C}^T$ 有拟逆函子 $\mathbb{L}$. 由于分裂对, 分裂叉, 生 $\varinjlim$ 等性质不受范畴等价影响, 而 $G = U^T \mathbb{K}$, 于是问题归结为证 $U^T: \mathcal{C}^T \to \mathcal{C}$ 对所有 $U^T$-分裂对生相应的余等化子, 这正是引理 7.6.12 的内容.
	
	其次有 (i) $\implies$ (iii). 既然已知 (i) $\implies$ (ii), 唯一任务是说明若 $G$ 是单子的, 则 $G$ 是保守的. 然而 $G = U^T \mathbb{K}$, 而 $\mathbb{K}$ 是范畴等价, $U^T$ 显然保守, 故 $G$ 也保守.
	
	以下证明 (ii) $\implies$ (i), 目标是构造 $\mathbb{K}$ 的拟逆函子 $\mathbb{L}$. 给定 $(M, a) \in \Obj(\mathcal{C}^T)$, 考虑 $\mathcal{D}$ 中的态射对
	$\begin{tikzcd}
		FGF(M) \arrow[shift left, r, "Fa"] \arrow[shift right, r, "\varepsilon_{FM}"'] & F(M).
	\end{tikzcd}$
	它们对 $G$ 的像
	$\begin{tikzcd}
		T^2(M) \arrow[shift left, r, "Ta"] \arrow[shift right, r, "\mu_M"'] & T(M)
	\end{tikzcd}$\!
	按例 7.6.10 扩充为分裂叉, 因此 $(Fa, \varepsilon_{FM})$ 是 $G$-分裂对, 从而原图可扩充为 $\mathcal{D}$ 中的余等化子图表, 亦即断言中的 (7.6.1).
	
	从 (7.6.1) 和余等化子的泛性质可见若有 $\mathcal{C}^T$ 的态射 $(M, a) \xrightarrow{\varphi} (M', a')$, 则有唯一态射 $\mathbb{L}(M, a) \to \mathbb{L}(M', a')$ 使下图交换:
	\begin{equation*}\begin{tikzcd}
		FGF(M) \arrow[shift left, r, "Fa"] \arrow[shift right, r, "\varepsilon_{FM}"'] \arrow[d, "FGF\varphi"'] & F(M) \arrow[r] \arrow[d, "F\varphi"] & \mathbb{L}(M, a) \arrow[d] \\
		FGF(M') \arrow[shift left, r, "{Fa'}"] \arrow[shift right, r, "\varepsilon_{FM'}"'] & F(M') \arrow[r] & \mathbb{L}(M', a'). 
	\end{tikzcd}\end{equation*}
	因此 $(M, a) \mapsto \mathbb{L}(M, a)$ 成为函子 $\mathbb{L}: \mathcal{C}^T \to \mathcal{D}$.
	
	兹证明 $\mathbb{L}\mathbb{K} \simeq \identity_{\mathcal{C}}$. 对所有 $N \in \Obj(\mathcal{D})$, 上述构造给出余等化子图表
	\begin{equation*}\begin{tikzcd}[column sep=large]
		FGFG(N) \arrow[shift left, r, "FG\varepsilon_N"] \arrow[shift right, r, "\varepsilon_{FG(N)}"'] & FG(N) \arrow[r] & \mathbb{L}\mathbb{K}(N).
	\end{tikzcd}\end{equation*}

	将此和引理 7.6.13 的余等化子图表相比较, 立得典范同构 $\mathbb{L}\mathbb{K} \simeq \identity_{\mathcal{C}}$.
	
	接着验证 $\mathbb{K}\mathbb{L} \simeq \identity_{\mathcal{C}^T}$. 易见 $\mathbb{K}F = \mathrm{Free}^T$, 两者都映对象 $M$ 为 $(GF(M), G\varepsilon_{FM})$. 按照定义和 $\mu := G\epsilon F$, 对 (7.6.1) 应用 $\mathbb{K}$ 的结果是
\renewcommand{\theequation}{7.6.2}
	\begin{equation}\label{eqn:Beck-aux3}
		\begin{tikzcd}[row sep=small]
			(T^2(M), \ldots) \arrow[shift left, r, "Ta"] \arrow[shift right, r, "\mu_M"'] & (T(M), \ldots) \arrow[r] & \mathbb{K}\mathbb{L}(M, a). \\
			\mathrm{Free}^T(T(M)) \arrow[equal, u] & \mathrm{Free}^T(M) \arrow[equal, u] &
		\end{tikzcd}
	\end{equation}

	对 (7.6.2) 左边两项运用忘却函子 $U^T$ 得到
	$\begin{tikzcd}
		T^2(M) \arrow[shift left, r, "{Ta}"] \arrow[shift right, r, "{\mu_M}"'] & T(M)
	\end{tikzcd}$,
	例 7.6.10 将之延长为以 $M$ 为余等化子的分裂叉. 然而 $U^T$ 生此余等化子 (引理 7.6.12), 回顾定义 1.5.2 可知 (7.6.2) 也是余等化子图表.
	
	另一方面, 将引理 7.6.13 施于伴随对 $(\mathrm{Free}^T, U^T)$ 和对象 $N = (M, a)$, 耐心展开此伴随对的详细刻画, 可得余等化子图表
	\[\begin{tikzcd}[row sep=tiny]
		\mathrm{Free}^T(T(M)) \arrow[shift left, r, "Ta"] \arrow[shift right, r, "\mu_M"'] & \mathrm{Free}^T(M) \arrow[r] & (M, a).
	\end{tikzcd}\]
	和上一段相比较立得 $\mathbb{K}\mathbb{L} \simeq \identity_{\mathcal{C}^T}$.
	
	最后说明 (iii) $\implies$ (ii). 注意到 (iii) 的后半部相当于说 $G$-分裂对皆有余等化子, 而且 $G$ 保此余等化子. 既然 $G$ 是保守函子, 注记 1.5.5 说明 (ii) 成立.
	
	若将一切态射倒转, 但函子走向不变, 亦即以 $\mathcal{C}^{\opp}$ 和 $\mathcal{D}^{\opp}$ 代 $\mathcal{C}$ 和 $\mathcal{D}$, 便可以得到余单子性的刻画. 此处不赘.
\end{proof}
\subsection*{必要作者背景与局部支撑（不另计题）}
\subsection*{作者标准术语：原书定义 1.5.2、1.5.4 与注记 1.5.5}
\noindent\textit{整理说明：以下原文从同一版本 PDF 物理页 45--46（书内页 37--38）转录；锥范畴采用通常定义。}
\par\medskip
\def\CorpusNumber{1.5.2}\begin{definition}\label{def:create-limit}
给定 $F: \mathcal{C}\to\mathcal{D}$ 和函子 $\alpha:I\to\mathcal{C}$, 考虑锥范畴之间的函子 $(\alpha/\Delta)\to(F\alpha/\Delta)$.
\begin{itemize}
\item 称 $F$ \textbf{保}此 $\varinjlim$, 如果 $(\alpha/\Delta)\to(F\alpha/\Delta)$ 映始对象 (假如存在) 为始对象.
\item 称 $F$ \textbf{返}此 $\varinjlim$, 如果仅 $(\alpha/\Delta)$ 的始对象方能被映为 $(F\alpha/\Delta)$ 的始对象.
\item 称 $F$ \textbf{生}此 $\varinjlim$, 如果一旦 $\varinjlim F\alpha$ 存在, 则 $\varinjlim\alpha$ 存在, 而且此时 $F$ 保此 $\varinjlim$ 也返此 $\varinjlim$.
\end{itemize}
对于 $\varprojlim$ 同样有相应的概念, 以 $\mathcal{C}^{\opp}$ 和 $\mathcal{D}^{\opp}$ 代 $\mathcal{C}$ 和 $\mathcal{D}$ 即可相互过渡.
\end{definition}
\def\CorpusNumber{1.5.4}\begin{definition}\label{def:conservative}
满足以下条件的函子 $F$ 称为\textbf{保守}的: 态射 $f\in\Mor(\mathcal{C})$ 为同构当且仅当 $Ff\in\Mor(\mathcal{D})$ 为同构.
\end{definition}
\def\CorpusNumber{1.5.5}\begin{remark}\label{rem:conservative-reflect}
若 $F$ 是保守函子, 则只要 $\varinjlim\alpha$ 存在, 而且 $F$ 保此 $\varinjlim$, 则 $F$ 也必然返此 $\varinjlim$. 缘由如下: 设 $L\in\Obj((\alpha/\Delta))$ 被映为 $(F\alpha/\Delta)$ 的始对象; 考虑 $(\alpha/\Delta)$ 中的唯一态射 $\varinjlim\alpha\to L$, 由于 $F$ 保此 $\varinjlim$, 它在 $F$ 之下的像必为同构, 从而保守条件确保 $\varinjlim\alpha\to L$ 也是同构.

在此前提下, 关于 $F$ 生 $\varinjlim$ 的定义中可以去除返 $\varinjlim$ 的条件.
\end{remark}
\par\medskip\noindent\textit{作者原文：chapter7.tex，行 1281--1369.}\par
设 $\mathcal{C}$ 是范畴. 考虑以所有自函子 $T: \mathcal{C} \to \mathcal{C}$ 为对象, 以自函子之间的态射构成的函子范畴 $\mathcal{C}^{\mathcal{C}}$. 它自然地成为一个幺半范畴, 其乘法 $\otimes$ 是函子的合成运算, 幺元 $\munit$ 是恒等函子. 因为函子的合成满足严格结合律, 它还是严格幺半范畴.

\def\CorpusNumber{7.6.1}
\begin{definition}[单子和余单子]\label{def:monad}
	\index{danzi@单子 (monad)}
	\index{yudanzi@余单子 (comonad)}
	设 $\mathcal{C}$ 为范畴. 幺半范畴 $\mathcal{C}^{\mathcal{C}}$ 中的代数 (定义 7.1.1) 称为 $\mathcal{C}$ 上的单子. 对偶地, $\mathcal{C}^{\opp}$ 上的单子称为 $\mathcal{C}$ 上的余单子.
\end{definition}

对自函子范畴 $\mathcal{C}^{\mathcal{C}}$ 展开定义 7.1.1, 可知单子是资料 $(T, \mu, \eta)$, 其中 $T: \mathcal{C} \to \mathcal{C}$ 是函子, $\mu: T^2 \to T$ 和 $\eta: \identity_{\mathcal{C}} \to T$ 是态射, 使得下图交换
\begin{equation*}\begin{gathered}
	\begin{tikzcd}
		T \arrow[r, "{\eta T}"] \arrow[rd, "\identity"'] & T^2 \arrow[d, "\mu"] & T \arrow[l, "{T\eta}"'] \arrow[ld, "\identity"] \\
		& T &
	\end{tikzcd} \quad
	\begin{tikzcd}
		T^3 \arrow[r, "\mu T"] \arrow[d, "T\mu"'] & T^2 \arrow[d, "\mu"] \\
		T^2 \arrow[r, "\mu"'] & T.
	\end{tikzcd}
\end{gathered}\end{equation*}

相较于原定义, 此处不再需要结合约束, 而 $\eta \otimes \identity$ (或 $\identity \otimes \eta$) 则被翻译为 $\eta T$ (或 $T\eta$), 依此类推.

对偶地, 余单子 $(L, \delta, \epsilon)$ 由函子 $L: \mathcal{C} \to \mathcal{C}$, 态射 $\delta: L \to L^2$ 和 $\epsilon: L \to \identity_{\mathcal{C}}$ 组成, 条件是使下图交换.
\begin{equation*}\begin{gathered}
	\begin{tikzcd}
		L  & L^2 \arrow[l, "{\epsilon L}"'] \arrow[r, "{L\epsilon}"] & L \\
		& L \arrow[lu, "\identity"] \arrow[ru, "\identity"'] \arrow[u, "\delta"] &
	\end{tikzcd} \quad
	\begin{tikzcd}
		L^3 & L^2 \arrow[l, "\delta L"'] \\
		L^2 \arrow[u, "L\delta"] & L \arrow[u, "\delta"'] \arrow[l, "\delta"]
	\end{tikzcd}
\end{gathered}\end{equation*}

我们习惯将资料 $(T, \mu, \eta)$ (或 $(L, \delta, \epsilon)$) 简记为 $T$ (或 $L$). 单子和余单子的主要来源是伴随对.

\def\CorpusNumber{7.6.2}
\begin{example}[伴随对确定单子]\label{eg:adjunction-monad}
	考虑一对伴随函子
	\[\begin{tikzcd}
		F: \mathcal{C} \arrow[shift left, r] & \mathcal{D}: G \arrow[shift left, l],
	\end{tikzcd}\]
	以及相应的单位 $\eta: \identity_{\mathcal{C}} \to GF$ 和余单位 $\varepsilon: FG \to \identity_{\mathcal{D}}$ 态射. 今将定义 $\mathcal{C}$ 上的单子 $(T, \mu, \eta)$ 和 $\mathcal{D}$ 上的余单子 $(L, \delta, \varepsilon)$ 如下.
	\[\begin{array}{|c|c|}\hline
		T := GF & L := FG \\
		\mu := \left[ GFGF \xrightarrow{G\varepsilon F} GF \right] & \delta := \left[ FG \xrightarrow{F\eta G} FGFG \right] \\
		\eta : \identity_{\mathcal{C}} \to GF & \varepsilon: FG \to \identity_{\mathcal{D}} \\ \hline 
	\end{array}\]

	基于对偶性, 就 $(T, \mu, \eta)$ 的情形验证单子的公理即可. 首先是关于幺元的图表
	\[\begin{tikzcd}[column sep=large]
		GF \arrow[r, "{\eta GF}"] \arrow[equal, rd] & GFGF \arrow[d, "{G\varepsilon F}"] & GF \arrow[l, "{GF\eta}"'] \arrow[equal, ld] \\
		& GF &
	\end{tikzcd}\]
	其交换性来自三角等式 $(G\varepsilon)(\eta G) = \identity_G$ 和 $(\varepsilon F)(F\eta) = \identity_F$. 结合律图表是
	\[\begin{tikzcd}[column sep=large]
		GFGFGF \arrow[d, "{G\varepsilon FGF}"'] \arrow[r, "{GFG\varepsilon F}"] & GFGF \arrow[d, "G\varepsilon F"] \\
		GFGF \arrow[r, "G\varepsilon F"'] & GF
	\end{tikzcd}\]
	它是交换的: 两路合成同样图解为
	\[\begin{tikzcd}
		\mathcal{C} \arrow[r, "F"] & \mathcal{D} \arrow[r, "G"] \arrow[rr, bend right=60, ""{name=A}, "{\identity}"'] & \mathcal{C} \arrow[r, "F"] \arrow[Rightarrow, to=A, "\varepsilon"] &
		\mathcal{D} \arrow[r, "G"] \arrow[rr, bend right=60, ""{name=B}, "{\identity}"'] & \mathcal{C} \arrow[r, "F"] \arrow[Rightarrow, to=B, "\varepsilon"] & \mathcal{D} \arrow[r, "G"] & \mathcal{C}
	\end{tikzcd}\]
	只是两个弓形区域的合成次序不同, 其产物则相同; 这是态射纵横合成的互换律 \cite[引理 2.2.7]{Li1} 的一则特例.
\end{example}

由于自函子作用在 $\mathcal{C}$ 上, 我们希望在 $\mathcal{C}$ 中探讨在单子 $T$ 作用下的``模''. 由于历史的原因, 这种结构也被称为 $T$-代数.

\def\CorpusNumber{7.6.3}
\begin{definition}[S.\ Eilenberg, J.\ C.\ Moore]\label{def:T-Alg}
	\index{mo}
	设 $(T, \mu, \eta)$ 是 $\mathcal{C}$ 上的单子. 所谓 $T$-模, 系指资料 $(M, a)$, 其中 $M \in \Obj(\mathcal{C})$ 而 $a$ 是 $\mathcal{C}$ 的态射 $T(M) \to M$, 使得下图交换
	\[\begin{tikzcd}
		T^2 (M) \arrow[r, "Ta"] \arrow[d, "\mu_M"'] & T(M) \arrow[d, "a"] \\
		T(M) \arrow[r, "a"'] & M
	\end{tikzcd} \quad \begin{tikzcd}
		M \arrow[r, "\eta_M"] \arrow[rd, "{\identity_M}"'] & T(M) \arrow[d, "a"] \\
		& M .
	\end{tikzcd}\]

	所有 $T$-模构成范畴 $\mathcal{C}^T$: 从 $(M, a)$ 到 $(M', a')$ 的态射定为 $\mathcal{C}$ 中使得下图交换的态射 $f: M \to M'$
	\[\begin{tikzcd}
		T(M) \arrow[r, "a"] \arrow[d, "Tf"'] & M \arrow[d, "f"] \\
		T(M')  \arrow[r, "{a'}"'] & M' .
	\end{tikzcd}\]

	对偶地, 设 $(L, \delta, \epsilon)$ 是 $\mathcal{C}$ 上的余单子, 所谓 $L$-余模是资料 $(N, b)$, 其中 $b$ 是态射 $N \to L(N)$, 条件是有与先前相对偶的交换图表. 我们将 $L$-余模范畴记为 $\mathcal{C}^L$.
\end{definition}

关于 $T$-模的两个交换图表应当分别被设想为 $T$ 作用下的结合律和幺元律. 今后的例子和性质主要针对单子及其上的模加以陈述, 余单子和余模的版本纯然是对偶的.
\par\medskip\noindent\textit{作者原文：chapter7.tex，行 1401--1439.}\par
回到一般理论. 我们有忘却函子 $U^T: \mathcal{C}^T \to \mathcal{C}$ 映对象 $(M, a)$ 为 $M$. 以下给出忘却的左伴随: 自由.

\def\CorpusNumber{7.6.6}
\begin{definition-proposition}[自由 $T$-模]\label{def:free-T-Alg}
	\index{mo!自由 (free)}
	设 $(T, \mu, \eta)$ 是 $\mathcal{C}$ 上的单子. 定义函子
	\[ \mathrm{Free}^T: \mathcal{C} \to \mathcal{C}^T , \quad
	\begin{array}{rl}
		(\text{对象}\; M) & \mapsto (TM, \mu_M: T^2 M \to TM), \\
		(\text{态射}\; f: M \to M') & \mapsto Tf: TM \to TM'.
	\end{array}\]
	这给出伴随对
	\[\begin{tikzcd}
		\mathrm{Free}^T : \mathcal{C} \arrow[shift left, r] & \mathcal{C}^T: U^T \arrow[shift left, l],
	\end{tikzcd}\]
	而此伴随对在 $\mathcal{C}$ 上确定的单子正是 $(T, \mu, \eta)$.
\end{definition-proposition}
\begin{proof}
	为了验证 $(TM, \mu_M)$ 是 $T$-模, 需要的是交换图表
	\[\begin{tikzcd}
		T^3 (M) \arrow[r, "T\mu_M"] \arrow[d, "\mu_{T(M)}"'] & T^2(M) \arrow[d, "\mu_M"] \\
		T^2 (M) \arrow[r, "\mu_M"'] & T(M)
	\end{tikzcd} \quad \begin{tikzcd}
		T(M) \arrow[r, "\eta_{T(M)}"] \arrow[rd, "\identity"'] & T^2(M) \arrow[d, "\mu_M"] \\
		& T(M)
	\end{tikzcd}\]
	然而这是将定义 7.6.1 的函子图表在对象 $M$ 上取值的结果. 函子性是清晰的.
	
	其次, 留意到 $U^T \mathrm{Free}^T = T$. 取 $T$ 自带的 $\eta: \identity_{\mathcal{C}} \to T$, 并且定义
	\[ \epsilon: \mathrm{Free}^T U^T \to \identity_{\mathcal{C}^T}, \quad
	\epsilon_{(M, a)}: (TM, \mu_M) \xrightarrow{a} (M, a) \in \Obj(\mathcal{C}^T). \]
	定义 7.6.3 的左图说明 $\epsilon_{(M, a)}$ 是 $\mathcal{C}^T$ 中的态射. 它的函子性同样清楚.
	
	例行的验证说明 $(\mathrm{Free}^T, U^T)$ 构成分别以 $\eta$ 和 $\epsilon$ 为单位和余单位的伴随对, 细节不赘. 此外
	\begin{equation*}\begin{aligned}
		\left(U^T \epsilon \mathrm{Free}^T\right)_M & = U^T \epsilon_{(TM, \mu_M)} \\
		& = \left[ \mu_M: T^2(M) \to T(M)\right],
	\end{aligned}\end{equation*}
	这便说明伴随对 $(\mathrm{Free}^T, U^T)$ 在 $\mathcal{C}$ 上确定的单子是 $(T, \mu, \eta)$.
\end{proof}
\par\medskip\noindent\textit{作者原文：chapter7.tex，行 1445--1476.}\par
\def\CorpusNumber{7.6.7}
\begin{lemma}\label{prop:T-comparison}
	考虑来自伴随对 $(F, G)$ 的单子 $T$ 和余单子 $L$. 设 $N \in \Obj(\mathcal{D})$, 则资料
	\[ M := G(N) \in \Obj(\mathcal{C}), \quad a: T(M) = GFG(N) \xrightarrow{G\varepsilon_N} G(N) = M \]
	给出 $(M, a) \in \Obj(\mathcal{C}^T)$. 由此得到函子 $\mathbb{K}: \mathcal{D} \to \mathcal{C}^T$, 满足 $G = U^T \mathbb{K}$.
	
	对偶地, $F$ 也有典范分解 $\mathcal{C} \xrightarrow{\mathbb{K}} \mathcal{D}^L \to \mathcal{D}$, 今后写作 $F = U^L \mathbb{K}$, 其中 $U^L$ 是忘却函子\footnote{采用相同符号 $\mathbb{K}$ 不至于引起太大的混淆, 因为本书不会同时考虑单子和余单子.}.
\end{lemma}
\begin{proof}
	只论 $T$ 的情形. 定义 7.6.3 的图表对此化为
	\[\begin{tikzcd}
		GFGFG(N) \arrow[r, "GFG\varepsilon" inner sep=0.5em] \arrow[d, "{G\varepsilon FG}"'] & GFG(N) \arrow[d, "G\varepsilon"] \\
		GFG(N) \arrow[r, "G\varepsilon"'] & G(N)
	\end{tikzcd} \quad \begin{tikzcd}
		G(N) \arrow[r, "\eta G"] \arrow[equal, rd] & GFG(N) \arrow[d, "{G\varepsilon}"] \\
		& G(N),
	\end{tikzcd} \]
	交换性的验证和例 7.6.2 如出一辙, 不必重复; $N \mapsto (M, a)$ 的函子性是自明的.
\end{proof}

对于给定伴随对 $(F, G)$, 应用函子是一个丢失结构的过程. 引理 7.6.7 将 $G$ (或 $F$) 拆成 $U^T \mathbb{K}$ (或 $U^L \mathbb{K}$). 我们想明白结构是否只被第二步的忘却函子 $U^T$ 丢失, 如果答案是肯定的, 过程中丢失的信息便能借助单子 $T$ (或余单子 $L$) 来重构. 这就启发了以下概念.

\def\CorpusNumber{7.6.8}
\begin{definition}\label{def:monadic}
	考虑一对伴随函子
	$\begin{tikzcd}
		F: \mathcal{C} \arrow[shift left, r] & \mathcal{D}: G \arrow[shift left, l]
	\end{tikzcd}$,
	以此定义 $\mathcal{C}$ 上的单子 $T$ 和 $\mathcal{D}$ 上的余单子 $L$.
	\begin{itemize}
		\item 若引理 7.6.7 的函子 $\mathbb{K}: \mathcal{D} \to \mathcal{C}^T$ 是范畴等价, 则称此伴随对是\emph{单子的}.
		\item 若对偶版本 $\mathbb{K}: \mathcal{C} \to \mathcal{D}^L$ 是范畴等价, 则称此伴随对是\emph{余单子}的.
	\end{itemize}
\end{definition}
\par\medskip\noindent\textit{作者原文：chapter7.tex，行 1482--1572.}\par
\def\CorpusNumber{7.6.9}
\begin{definition-proposition}
	\index{fenliecha@分裂叉 (split fork)}
	考虑任意范畴 $\mathcal{C}$ 中的图表
	\begin{equation*}\label{eqn:split-coeq}\begin{tikzcd}[column sep=large]
		A \arrow[shift left, r, "u"] \arrow[shift right, r, "v" description] & B \arrow[r, "h"] \arrow[dashed, l, bend left, "t"] & Z \arrow[dashed, l, bend left, "s"]
	\end{tikzcd}\end{equation*}
	使得实线部分交换, 亦即 $hu = hv$, 而且 $h$ 和 $v$ 各自有虚线所示的截面 (亦即右逆) $s$ 和 $t$, 满足 $sh = ut \in \End(B)$. 具此性质的实线部分图表称为\emph{分裂叉}. 它自动给出 $u$ 和 $v$ 在 $\mathcal{C}$ 中的余等化子.
\end{definition-proposition}
\begin{proof}
	考虑以下场景
	\[\begin{tikzcd}[row sep=small]
		& & W \\
		A \arrow[shift left, r, "u"] \arrow[shift right, r, "v"'] & B \arrow[r, "h"'] \arrow[ru, "k"] & Z
	\end{tikzcd} \qquad \begin{array}{ll}
		W \in \Obj(\mathcal{C}) \\
		ku = kv.
	\end{array} \]
	若存在 $\varphi: Z \to W$ 使得 $\varphi h = k$, 则 $\varphi = \varphi hs = ks$; 反之, $\varphi := ks$ 确实满足 $\varphi h = ksh = kut = kvt = k$. 这就验证了泛性质.
\end{proof}

不同于一般的余等化子, 分裂叉是``绝对''的: 它们对任意函子的像仍是分裂叉.

\def\CorpusNumber{7.6.10}
\begin{example}\label{eg:T-split-fork}
	设 $(T, \mu, \eta)$ 是 $\mathcal{C}$ 上的单子, 则任何 $(M, a) \in \Obj(\mathcal{C}^T)$ 都给出 $\mathcal{C}$ 中的分裂叉
	\begin{equation*}\begin{tikzcd}[column sep=large]
			T^2(M) \arrow[shift left, r, "Ta"] \arrow[shift right, r, "\mu_M" description] & T(M) \arrow[r, "a"] \arrow[dashed, bend left, l, "\eta_{TM}"] & M \arrow[dashed, bend left, l, "\eta_M"]
		\end{tikzcd} \quad \begin{array}{ll}
			a (Ta) = a \mu_M , & \eta_M a = (Ta) \eta_{TM}, \\
			a \eta_M = \identity_M , & \mu_M \eta_{TM} = \identity_{TM}.
	\end{array}\end{equation*}
	右侧列出的等式是定义 7.6.1 和 7.6.3 的内容; 例如 $\eta_M a = (Ta) \eta_{TM}$ 缘于 $\eta: \identity_{\mathcal{C}} \to T$ 的自然性. 特别地, 上图将 $M$ 实现为 $Ta$ 和 $\mu_M$ 的余等化子.
\end{example}

\def\CorpusNumber{7.6.11}
\begin{definition}\label{def:Beck-split-pair}
	考虑函子 $G: \mathcal{D}\to \mathcal{C}$ 和 $\mathcal{D}$ 的一对态射 $f, g: X \rightrightarrows Y$. 如果存在态射 $h: G(Y) \to Z$ 使得图表
	\begin{equation*}\begin{tikzcd}
			G(X) \arrow[shift left, r, "Gf"] \arrow[shift right, r, "Gg"'] & G(Y) \arrow[r, "h"] & Z
	\end{tikzcd}\end{equation*}
	是 $\mathcal{C}$ 中的分裂叉, 则称 $(f, g)$ 是 $\mathcal{D}$ 中的 $G$-分裂对.
\end{definition}

设 $(f, g)$ 是 $G$-分裂对. 既然 $(Gf, Gg)$ 有余等化子, 一个自然的问题是能否将之提升为 $(f, g)$ 在 $\mathcal{D}$ 中的等化子, 以及此提升是否唯一. 余等化子是 $\varinjlim$ 的特例, 定义 1.5.2 关于函子生 $\varinjlim$ 或保 $\varinjlim$ 的概念为此提供了一套方便的术语.

留意到忘却函子 $U^T: \mathcal{C}^T \to \mathcal{C}$ 是定义 1.5.4 所谓的保守函子.

\def\CorpusNumber{7.6.12}
\begin{lemma}\label{prop:UT-split-fork}
	设 $(T, \mu, \eta)$ 是 $\mathcal{C}$ 上的单子, 则函子 $U^T$ 生 $U^T$-分裂对的余等化子.
\end{lemma}
\begin{proof}
	设有 $\mathcal{C}^T$ 中的一对态射 $u, v: (M, a) \rightrightarrows (M', a')$ 和 $\mathcal{C}$ 中的分裂叉
	\begin{equation*}\begin{tikzcd}[column sep=large]
		M \arrow[shift left, r, "u"] \arrow[shift right, r, "v" description] & M' \arrow[r, "h"] \arrow[dashed, l, bend left, "t"] & M'' \arrow[dashed, l, bend left, "s"].
	\end{tikzcd}\end{equation*}
	由于上图给出 $\mathcal{C}$ 中的等化子, 问题归结为将 $M''$ 唯一地扩充为 $(M'', a'') \in \Obj(\mathcal{C}^T)$, 使得 $h$ 是 $\mathcal{C}^T$ 中的态射, 然后说明这给出 $u$ 和 $v$ 在 $\mathcal{C}^T$ 中的余等化子.
	
	分裂叉对任意函子的像仍是分裂叉, 故在实线部分的交换图表
	\[\begin{tikzcd}
		T(M) \arrow[shift left, r, "Tu"] \arrow[shift right, r, "Tv"'] \arrow[d, "a"'] & T(M') \arrow[r, "Th"] \arrow[d, "{a'}"] & T(M'') \arrow[dashed, d, "{a''}"] \\
		M \arrow[shift left, r, "u"] \arrow[shift right, r, "v"'] & M' \arrow[r, "h"'] & M''
	\end{tikzcd}\]
	中, 两行都是 $\mathcal{C}$ 中的余等化子, 从而泛性质给出唯一的 $a''$ 使得全图交换. 如能说明 $(M'', a'')$ 是 $T$-代数, 则上图右部交换相当于说 $h$ 是 $\mathcal{C}^T$ 中的态射.
	
	为此, 问题在于验证以下交换图表
	\[\begin{tikzcd}
		M'' \arrow[r, "{\eta_{M''}}"] \arrow[equal, rd] & T(M'') \arrow[d, "{a''}"] \\
		& T(M'')
	\end{tikzcd} \quad
	\begin{tikzcd}
		T^2(M'') \arrow[r, "{\mu_{M''}}"] \arrow[d, "{Ta''}"'] & T(M'') \arrow[d, "{a''}"] \\
		T(M'') \arrow[r, "{a''}"'] & M'' .
	\end{tikzcd}\]
	余等化子图表说明 $h$ 和 $Th$ 满 (同理, $T^2 h$ 满), 故易从 $M$ 和 $M'$ 的相应交换图表说明上图确实交换.
	
	最后来说明 $h$ 给出 $u$ 和 $v$ 在 $\mathcal{C}^T$ 中的余等化子. 设有 $\mathcal{C}^T$ 的态射 $k: (M'', a'') \to (W, b)$ 满足 $ku = kv$, 则在 $\mathcal{C}$ 中存在唯一态射 $\varphi: M'' \to W$ 使得 $\varphi h = k$. 问题归结为证明 $\varphi$ 实则是 $\mathcal{C}^T$ 的态射. 考虑图表:
	\[\begin{tikzcd}
		T(M') \arrow[r, "Th"] \arrow[d, "{a'}"'] & T(M'') \arrow[r, "T\varphi"] \arrow[d, "{a''}"] & T(W) \arrow[d, "b"] \\
		M' \arrow[r, "h"'] & M'' \arrow[r, "\varphi"'] & W
	\end{tikzcd}\]
	上下两行分别合成为 $Tk$ 和 $k$, 因为 $h$ 和 $k$ 是 $\mathcal{C}^T$ 中的态射, 图表的整个外框和左方块皆交换, 故右侧方块和 $Th$ 合成以后交换. 已知 $Th$ 满, 故右侧方块本身交换, 这就说明了 $\varphi$ 是 $\mathcal{C}^T$ 中的态射.
\end{proof}

\def\CorpusNumber{7.6.13}
\begin{lemma}\label{prop:Beck-FG-coeq}
	设函子 $G: \mathcal{D} \to \mathcal{C}$ 有左伴随 $F$, 而且 $G$ 对所有 $G$-分裂对 $f, g: X \rightrightarrows Y$ (定义 7.6.11) 都生相应的余等化子, 则下图给出余等化子:
	\begin{equation*}\begin{tikzcd}[column sep=large]
		FGFG(N) \arrow[shift left, r, "FG\varepsilon_N"] \arrow[shift right, r, "\varepsilon_{FG(N)}"'] & FG(N) \arrow[r, "\varepsilon_N"] & N,
	\end{tikzcd}\end{equation*}
	其中 $N \in \Obj(\mathcal{D})$.
\end{lemma}
\begin{proof}
	对图表取 $G$ 给出例 7.6.10 在 $(M, a) := (G(N), G\varepsilon_N) = \mathbb{K}(N)$ 情形的分裂叉, 而 $G$ 生相应的余等化子. 故原图是余等化子.
\end{proof}
\subsection*{标准先修与保留说明（整理者）}
通常的伴随、锥与（余）极限、Abel 范畴、投射对象和生成元定义为标准先修。原书编号引用在本文件中保持原编号。森田分类中 Eilenberg--Watts 定理 7.7.2 为具名标准先修；本文件保留其陈述，以及实际用到的有限投射对偶、评价与余评价构造。原书命题 2.8.15 为投射生成元的忠实正合刻画；A.2.2 为紧对象定义（已在局部背景中复述）。原文的简短例行验证及任何已存在的笔误均原样保留，未作数学修改。
\begin{thebibliography}{Li1}
\bibitem{Li1} 李文威，《代数学方法》第一卷：基础架构，高等教育出版社。原书引用保留原章节和结果编号；\url{https://www.wwli.asia/mathematics/AlJabr/}.
\end{thebibliography}
\end{document}
